\section{Two informative raw-kernel realizations}
\label{sec:realizations}
The protocol is the gate--acquisition instrument of Theorem~\ref{thm:gate}. Let $Y\sim\nu_K$ and $T$ have distribution $\Pp(T\le s)=s^\alpha$, independently of each other and of the gate coins. The primary active report is $(Y,T)$, with $f(Y,T)=Y$ and $w(Y,T)=T^2$. Choose one coordinate vector $e_1$ and a compact convex $H$ containing $K+[-1,1]e_1$. The complete report includes an informative sign $V$. The raw driver projection is common; the full $V$ likelihood is not required to be common or independent of the physical state.

\begin{theorem}[Controlled hidden-state realization]\label{thm:classical}
At every raw call a hidden sign $E$ first evolves by an action-selected transition matrix whose entries belong to $[q,1-q]$, with fixed $0<q\le1/2$. The matrices may depend on an unknown fixed model. A gate call applies its independent Bernoulli protocol coin and permits no stop. At an active call, draw $(Y,T)$ as specified, then report
\[
 \Pp(V=v\mid E=e)=\frac{1+\eta ve}{2},\qquad 0<\eta<1.
\]
After commitment audit $Z=Y+TEe_1$ without a new transition. Rejection keeps the hidden sign and returns the visible protocol to the gate. Then Theorems~\ref{thm:transfer} and \ref{thm:gate} hold uniformly over every admissible transition family and every $0<p\le1$. The law remains nonregenerative when its transitions have distinct rows.
\end{theorem}
\begin{proof}
All operations are normalized positive kernels. Before an active report, either sign has conditional probability at least $q$, for every surviving history and every chosen action. Let $\gamma=(1-\eta)/2$ and $a=q\gamma/(1-\gamma)>0$. After observing $V$, both posterior sign probabilities are at least $a$. If $b$ is the posterior sign mean, then
\[
 \mu=Y+Tb e_1,\qquad v=T^2(1-b^2),\qquad
 4a(1-a)T^2\le v\le T^2,\quad \|\mu-Y\|^2\le T^2.
\]
Thus Definition~\ref{def:calibration} holds with $B_0=0$, $a_v=4a(1-a)$ and $A_v=a_b=1$. These inequalities hold conditional on $V$, so the all-history lower allows stopping and action selection informed by it. The exact primary projection follows independently by marginalizing the raw kernel over the sign and $V$; gate and active laws are the displayed probabilities, at every surviving history. The true prior generally changes through rejections, and its effect on the posterior $b$ does not disappear from the physical experiment.

For each fixed model, the finite sign instrument acts linearly on the two unnormalized sign masses. In the unknown-model interpretation these carriers are kept as a model-indexed family, or as a joint posterior only after a prior is specified; two masses alone do not encode uncertainty about a transition matrix. Include the reported $Y,T$, the protocol mode and legal action signature for a possible immediate audit. Report probabilities and the audit moments are executable functions of this carrier. Distinct beliefs are separated by future informative observations whenever a separating transition is admitted; degenerate directions are quotiented out, not declared identifiable. Report total variation is Lipschitz in the input sign law. Normalized active updates have denominators at least $\gamma$ and are uniformly continuous on the compact belief interval. The observer does not store or numerically iterate this belief. Its exact stored protocol mode and the killed-flow realization supply the causal implementation. Proposition~\ref{prop:geometry} verifies the global cover and mass bounds from the actual $(Y,T)$ law. Constants depend on $q,\eta$ but not on the gate rate $p$.
\end{proof}

\begin{theorem}[Noncommuting quantum realization]\label{thm:quantum}
Prepare a qubit in $I/2$. At each raw call apply the positivity-improving channel
\[
 D_q(\rho)=(1-2q)\rho+q\,\tr(\rho)I,\qquad0<q\le1/2,
\]
followed by an action-selected unitary $U_a$. A gate call applies only the independent Bernoulli protocol coin and has no terminal readout. At an active call independently report $(Y,T)$, use effects
\[
 K_{a,v}=\left(\frac{I+\eta v\,n_a\cdot\sigma}{2}\right)^{1/2},\qquad
 \mathcal I_{a,v}(\rho)=K_{a,v}U_aD_q(\rho)U_a^*K_{a,v},
\]
and report $v$. At commitment measure $\sigma_z$ on the normalized current state, obtaining $B\in\{-1,1\}$, and audit $Z=Y+TB e_1$. Rejection retains the quantum state and returns to the gate. Nonparallel directions $n_a$ are allowed. The same two general theorems hold, with the same dependence on $p,N,M,d,\alpha$ as in the classical realization.
\end{theorem}
\begin{proof}
The channel and each instrument operation are completely positive, and $\sum_vK_{a,v}^2=I$. Gate and active primary variables are physically sampled independently, giving \eqref{eq:projection}. Put $\gamma=(1-\eta)/2$. Before a weak measurement, $U_aD_q(\rho)U_a^*\succeq qI$. Every effect has eigenvalues in $[\gamma,1-\gamma]$, so its probability is in that interval and its normalized output obeys
\[
 \rho'\succeq\frac{q\gamma}{1-\gamma}I=aI.
\]
The subsequent Pauli mean has absolute value at most $1-2a$. Its variance is therefore at least $4a(1-a)$ and at most one. Scaling by $T$ and shifting by $Y$ gives exactly the calibration inequalities proved above, now derived from the quantum raw kernel and Born audit, not imported from the classical theorem.

For each fixed model, positive operations on finite Hermitian matrices and executable Born effects give an unnormalized predictive representation. A uniform unknown-model quotient retains the model-indexed family, not a density matrix that silently encodes the true model. A tomographically separating probe family, when admitted, identifies density operators on the attained image; otherwise the corresponding redundant directions are factored out. The observable protocol signature is included in either case. Report probabilities are Lipschitz in trace norm, and normalization has denominator at least $\gamma$, proving report and predictive-update continuity. No infinite-dimensional operator closure is inferred. For $q<1/2$, invertible weak effects can retain coherence after a rejection. Choosing axes $x$ and $z$ gives noncommuting effects. The physical qubit is not an available observer memory register: the implementation has precisely the charged classical live and terminal labels. Finally the actual primary location/reliability distribution verifies \eqref{eq:rectangles} independently of the quantum posterior.
\end{proof}

Weak measurement and its conditioning interpretation are standard~\cite{Jacobs}. The added protocol describes heralded preparation attempts followed by a measurement; every failure and subsequent attempt is paid. The examples deliberately verify a sufficient calibration class. The estimates do not become uniform as the hidden positivity floor $q$ vanishes; their new uniformity is in the independently varying physical waiting rate $p$.
